\section{Embedded points}
\label{section-embedded}

\noindent
Let $R$ be a ring and let $M$ be an $R$-module.
Recall that a prime $\mathfrak p \subset R$ is an
{\it embedded associated prime} of $M$ if it is an associated prime of
$M$ which is not minimal among the associated primes of $M$. See
Algebra, Definition \ref{algebra-definition-embedded-primes}.
Here is the definition of embedded associated points
for quasi-coherent sheaves on schemes
as given in \cite[IV Definition 3.1.1]{EGA}.

\begin{definition}
\label{definition-embedded}
Let $X$ be a scheme.
Let $\mathcal{F}$ be a quasi-coherent sheaf on $X$.
\begin{enumerate}
\item An {\it embedded associated point} of $\mathcal{F}$
is an associated point which is not maximal among the
associated points of $\mathcal{F}$, i.e., it is the specialization
of another associated point of $\mathcal{F}$.
\item A point $x$ of $X$ is called an {\it embedded point}
if $x$ is an embedded associated point of $\mathcal{O}_X$.
\item An {\it embedded component} of $X$ is an irreducible
closed subset $Z = \overline{\{x\}}$ where $x$ is an embedded
point of $X$.
\end{enumerate}
\end{definition}

\noindent
In the Noetherian case when $\mathcal{F}$ is coherent we have
the following.

\begin{lemma}
\label{lemma-embedded}
Let $X$ be a locally Noetherian scheme.
Let $\mathcal{F}$ be a coherent $\mathcal{O}_X$-module.
Then
\begin{enumerate}
\item the generic points of irreducible components of
$\text{Supp}(\mathcal{F})$ are associated points of $\mathcal{F}$, and
\item an associated point of $\mathcal{F}$ is embedded if and only
if it is not a generic point of an irreducible component
of $\text{Supp}(\mathcal{F})$.
\end{enumerate}
In particular an embedded point of $X$ is an associated point of $X$
which is not a generic point of an irreducible component of $X$.
\end{lemma}

\begin{proof}
Recall that in this case $Z = \text{Supp}(\mathcal{F})$ is closed, see
Morphisms, Lemma \ref{morphisms-lemma-support-finite-type}
and that the generic points of irreducible components of $Z$ are
associated points of $\mathcal{F}$, see
Lemma \ref{lemma-minimal-support-in-ass}.
Finally, we have $\text{Ass}(\mathcal{F}) \subset Z$, by
Lemma \ref{lemma-ass-support}.
These results, combined with the fact that $Z$ is a sober topological
space and hence every point of $Z$ is a specialization of a generic
point of $Z$, imply (1) and (2).
\end{proof}

\begin{lemma}
\label{lemma-S1-no-embedded}
Let $X$ be a locally Noetherian scheme.
Let $\mathcal{F}$ be a coherent sheaf on $X$.
Then the following are equivalent:
\begin{enumerate}
\item $\mathcal{F}$ has no embedded associated points, and
\item $\mathcal{F}$ has property $(S_1)$.
\end{enumerate}
\end{lemma}

\begin{proof}
This is Algebra, Lemma \ref{algebra-lemma-criterion-no-embedded-primes},
combined with Lemma \ref{lemma-associated-affine-open} above.
\end{proof}

\begin{lemma}
\label{lemma-noetherian-dim-1-CM-no-embedded-points}
Let $X$ be a locally Noetherian scheme of dimension $\leq 1$.
The following are equivalent
\begin{enumerate}
\item $X$ is Cohen-Macaulay, and
\item $X$ has no embedded points.
\end{enumerate}
\end{lemma}

\begin{proof}
Follows from Lemma \ref{lemma-S1-no-embedded} and the definitions.
\end{proof}

\begin{lemma}
\label{lemma-scheme-theoretically-dense-contain-embedded-points}
Let $X$ be a locally Noetherian scheme. Let $U \subset X$ be an
open subscheme. The following are equivalent
\begin{enumerate}
\item $U$ is scheme theoretically dense in $X$
(Morphisms, Definition \ref{morphisms-definition-scheme-theoretically-dense}),
\item $U$ is dense in $X$ and $U$ contains all embedded points of $X$, and
\item $U$ contains all associated points of $X$.
\end{enumerate}
\end{lemma}

\begin{proof}
The question is local on $X$, hence we may assume that $X = \Spec(A)$
where $A$ is a Noetherian ring. Then $U$ is quasi-compact
(Properties, Lemma \ref{properties-lemma-immersion-into-noetherian})
hence $U = D(f_1) \cup \ldots \cup D(f_n)$
(Algebra, Lemma \ref{algebra-lemma-qc-open}).
In this situation $U$ is scheme theoretically dense in $X$ if and only if
$A \to A_{f_1} \times \ldots \times A_{f_n}$ is injective, see
Morphisms, Example \ref{morphisms-example-scheme-theoretic-closure}.
Condition (3) translated into algebra means that for every associated
prime $\mathfrak p$ of $A$ there exists an $i$ with $f_i \not \in \mathfrak p$.

\medskip\noindent
Assume (1), i.e., $A \to A_{f_1} \times \ldots \times A_{f_n}$ is injective.
Say $a \in A$ has annihilator a prime $\mathfrak p$. Since $a$ maps
to a nonzero element of $A_{f_i}$ for some $i$ we see that
$f_i \not \in \mathfrak p$. Thus (3) holds.

\medskip\noindent
Note that $U$ is dense in $\Spec(A)$ if and only if it contains all
minimal prime ideals. Since the set of associated primes
of $A$ is equal to the union of the set of minimal primes and the
set of embedded associated primes (see Algebra, Proposition
\ref{algebra-proposition-minimal-primes-associated-primes}),
we see that (2) and (3) are equivalent.

\medskip\noindent
Assume (3). This means that every associated prime $\mathfrak p$ of $A$
corresponds to a prime of $A_{f_i}$ for some $i$. Then
$A \to A_{f_1} \times \ldots \times A_{f_n}$ is injective because
$A \to \prod_{\mathfrak p \in \text{Ass}(A)} A_\mathfrak p$ is injective
by Algebra, Lemma \ref{algebra-lemma-zero-at-ass-zero}. In this way
we see that (3) implies (1).
\end{proof}

\begin{lemma}
\label{lemma-remove-embedded-points}
Let $X$ be a locally Noetherian scheme.
Let $\mathcal{F}$ be a coherent sheaf on $X$.
The set of coherent subsheaves
$$
\{
\mathcal{K} \subset \mathcal{F}
\mid
\text{Supp}(\mathcal{K})\text{ is nowhere dense in }\text{Supp}(\mathcal{F})
\}
$$
has a maximal element $\mathcal{K}$.
Setting $\mathcal{F}' = \mathcal{F}/\mathcal{K}$ we have the
following
\begin{enumerate}
\item $\text{Supp}(\mathcal{F}') = \text{Supp}(\mathcal{F})$,
\item $\mathcal{F}'$ has no embedded associated points, and
\item there exists a dense open $U \subset X$ such that
$U \cap \text{Supp}(\mathcal{F})$ is dense in $\text{Supp}(\mathcal{F})$
and $\mathcal{F}'|_U \cong \mathcal{F}|_U$.
\end{enumerate}
\end{lemma}

\begin{proof}
This follows from
Algebra, Lemmas \ref{algebra-lemma-remove-embedded-primes} and
\ref{algebra-lemma-remove-embedded-primes-localize}.
Note that $U$ can be taken as the complement of the closure
of the set of embedded associated points of $\mathcal{F}$.
\end{proof}

\begin{lemma}
\label{lemma-no-embedded-points-endos}
Let $X$ be a locally Noetherian scheme.
Let $\mathcal{F}$ be a coherent $\mathcal{O}_X$-module
without embedded associated points. Set
$$
\mathcal{I}
=
\Ker(\mathcal{O}_X
\longrightarrow
\SheafHom_{\mathcal{O}_X}(\mathcal{F}, \mathcal{F})).
$$
This is a coherent sheaf of ideals which defines a closed
subscheme $Z \subset X$ without embedded points. Moreover
there exists a coherent sheaf $\mathcal{G}$ on $Z$
such that (a) $\mathcal{F} = (Z \to X)_*\mathcal{G}$,
(b) $\mathcal{G}$ has no associated embedded points, and
(c) $\text{Supp}(\mathcal{G}) = Z$ (as sets).
\end{lemma}

\begin{proof}
Some of the statements we have seen in the proof of
Cohomology of Schemes, Lemma \ref{coherent-lemma-coherent-support-closed}.
The others follow from
Algebra, Lemma \ref{algebra-lemma-no-embedded-primes-endos}.
\end{proof}



